\chapter*{Conclusion et pistes d'approfondissement}
\addcontentsline{toc}{chapter}{Conclusion et pistes d'approfondissement}

Ce document a présenté les deux grandes approches actuelles de génération d'images par transport
continu de distribution — diffusion et flow matching — depuis leurs fondements mathématiques jusqu'à
une implémentation complète et fonctionnelle. Les deux méthodes partagent bien plus qu'il n'y paraît
au premier abord : la même intuition (remplacer un grand saut génératif par une longue séquence de
petits pas), la même architecture de réseau, et ne diffèrent fondamentalement que par la façon dont le
chemin entre bruit et donnée est construit (stochastique et non-linéaire pour le DDPM, déterministe et
linéaire pour le flow matching).

\paragraph{Pistes d'approfondissement.} Plusieurs directions naturelles prolongent ce travail :
\begin{itemize}
    \item \textbf{Diffusion latente.} Remplacer la diffusion pixel-space par une diffusion dans
    l'espace latent d'un VAE (comme Stable Diffusion), pour viser une résolution supérieure à
    ressources de calcul égales.
    \item \textbf{Schémas d'échantillonnage accélérés} pour le DDPM (DDIM, solveurs d'ordre supérieur),
    qui réduisent le nombre d'étapes de génération sans réentraîner le réseau.
    \item \textbf{Fine-tuning de poids pré-entraînés} : charger des poids issus d'un modèle de
    diffusion existant dans l'architecture U-Net de ce projet, plutôt que d'entraîner entièrement
    \emph{from scratch}.
    \item \textbf{Évaluation quantitative} (FID, CLIP-score) pour comparer objectivement les deux
    approches au-delà de l'inspection visuelle.
\end{itemize}

\begin{thebibliography}{9}

\bibitem{ho2020ddpm}
J. Ho, A. Jain, P. Abbeel.
\textit{Denoising Diffusion Probabilistic Models}.
NeurIPS, 2020.

\bibitem{ho2022cfg}
J. Ho, T. Salimans.
\textit{Classifier-Free Diffusion Guidance}.
NeurIPS Workshop on Deep Generative Models, 2022.

\bibitem{lipman2023flow}
Y. Lipman, R. T. Q. Chen, H. Ben-Hamu, M. Nickel, M. Le.
\textit{Flow Matching for Generative Modeling}.
ICLR, 2023.

\bibitem{radford2021clip}
A. Radford, J. W. Kim, C. Hallacy, et al.
\textit{Learning Transferable Visual Models From Natural Language Supervision (CLIP)}.
ICML, 2021.

\bibitem{jiang2021celebadialog}
Y. Jiang, Z. Huang, X. Pan, C. C. Loy, Z. Liu.
\textit{Talk-to-Edit: Fine-Grained Facial Editing via Dialog}.
ICCV, 2021.

\bibitem{rombach2022ldm}
R. Rombach, A. Blattmann, D. Lorenz, P. Esser, B. Ommer.
\textit{High-Resolution Image Synthesis with Latent Diffusion Models}.
CVPR, 2022.

\bibitem{ronneberger2015unet}
O. Ronneberger, P. Fischer, T. Brox.
\textit{U-Net: Convolutional Networks for Biomedical Image Segmentation}.
MICCAI, 2015.

\bibitem{vaswani2017attention}
A. Vaswani, N. Shazeer, N. Parmar, et al.
\textit{Attention Is All You Need}.
NeurIPS, 2017.

\end{thebibliography}
